\section{Conclusion}\label{sec:r31conclusion}
Neural Bellman Operators separate evaluation of an economic future from the
construction of feasible decisions. The paper connects this architecture to
economic accuracy at three levels: centered decision errors, full-policy
comparison, and constructive training of a specified neural continuation.
The primitive-budget theorem determines a sufficient learning schedule before
its own-policy targets exist and propagates the resulting approximation
error through a nonlinear recursive preference criterion. Residual-centered
outward evaluation then assesses the rounded returned policy, including its
explicit execution-error allowance.

The full-policy economy has continuous Gaussian uncertainty, vector actions,
and reoptimized controls at every date. Its certificates cover all real
states and all adapted finite-recursive-cost deviations. These properties
are different from a finite-catalogue scalar-decision certificate, and the
evidence keeps that distinction explicit. The construction uses the candidate's
own finalized future rather than optimal-policy training labels.

Numerical necessity remains a separate comparison. Analytically structured
capital economies admit inexpensive conventional and structural solutions;
those solutions are retained and charged on the same basis as NBO. The
reported certificates and finite training bounds do not turn additional
neural fitting work into an unobserved computational advantage. Earlier
simulation, refresh, mechanism, and application results remain part of the
record, including their unfavorable comparisons. The contribution is the
explicit link between a trainable Bellman representation, its economic
policy target, and the work required to construct and verify it in the
stated classes. Its extension to a different utility domain, nonlinear
technology, or strategic equilibrium requires that model's own approximation,
domain, and deviation accounts, rather than a change in the meaning of the
reported accuracy.
